\documentclass{article}
\usepackage{amsmath, amssymb, amsthm}
\usepackage{hyperref}
\usepackage{geometry}
\geometry{margin=1in}

\title{Erd\H{o}s Problem \#611}
\date{Last updated: 2025-08-31}
\author{Source: erdosproblems.com}

\begin{document}
\maketitle

\noindent\textbf{Status:} OPEN \quad \textbf{No prize}

\noindent\textbf{Tags:} graph theory

\medskip

\noindent\textbf{Problem Statement:}

For a graph $G$ let $\tau(G)$ denote the minimal number of vertices that include at least one from each maximal clique of $G$ (sometimes called the clique transversal number). Is it true that if all maximal cliques in $G$ have at least $cn$ vertices then $\tau(G)=o_c(n)$? Similarly, estimate for $c&#62;0$ the minimal $k_c(n)$ such that if every maximal clique in $G$ has at least $k_c(n)$ vertices then $\tau(G)&#60;(1-c)n$.




A problem of Erd\H{o}s, Gallai, and Tuza \cite{EGT92}, who proved for the latter question that $k_c(n) \geq n^{c'/\log\log n}$ for some $c'&gt;0$, and that if every clique has size least $k$ then $\tau(G) \leq n-(kn)^{1/2}$. Bollob\'{a}s and Erd\H{o}s proved that if every maximal clique has at least $n+3-2\sqrt{n}$ vertices then $\tau(G)=1$ (and this threshold is best possible). See also [610] and the entry in the graphs problem collection .




References

[EGT92] Erd\H{o}s, Paul and Gallai, Tibor and Tuza, Zsolt, Covering the cliques of a graph with vertices . Discrete Math. (1992), 279-289.


\bigskip
\noindent\small{Source: \url{https://www.erdosproblems.com/611}}
\end{document}
